\part{Nonclassical Logic}
\label{part:nonclassical}

\chapter{Intuitionistic Logic}
\label{ch:intuitionistic-logic}

\begin{convention}[Nonclassical formula metavariables]
Throughout this part, \(A,B,C\) range over formulas and \(p,q\) range over
propositional variables.  The symbols \(\Gamma,\Delta,\Sigma,\Pi\) retain their
normative standing as formula contexts.  Any narrower local role is declared
before use.
\end{convention}

Intuitionistic logic is specified positively by its calculus and semantics.  It
is not defined as classical logic with selected theorems deleted.  This chapter
uses \(\LJ\) from \cref{ch:lk-lj} as the calculus and Kripke
structures for its semantics; see
\cite{Kripke1965,TroelstraSchwichtenberg2000}.

